\documentclass[11pt]{article}
\usepackage{amsmath,amssymb}
\usepackage{physics}
\usepackage[margin=1in]{geometry}
\usepackage{enumitem}
\usepackage{qcircuit}

\title{ATPL 2025 --- Exercises: State Preparation and the Input Problem}
\author{Lecture 1: Why do we need hybrid programming models?}
\date{}

\begin{document}
\maketitle

Background: the "Input Problem" slides from Lecture 1. For the standard
recursive state-preparation construction referenced below, see e.g.\
M{\o}tt{\"o}nen, Vartiainen, Bergholm \& Salomaa, \emph{Transformation of
quantum states using uniformly controlled rotations} (2004), or Shende,
Bullock \& Markov (2006).

\begin{enumerate}[label=\textbf{Exercise \arabic*.}, leftmargin=2.2cm]


%%%%%%%%%%%%%%%%%%%%%%%%%%%%%%%%%%%%
%
%  Basis-state preparation is cheap
%
%%%%%%%%%%%%%%%%%%%%%%%%%%%%%%%%%%%%

\item \textbf{Basis-state preparation is cheap.}\\
Given a target computational basis state $\ket{b_1 b_2 \cdots b_n}$ with
$b_k \in \{0,1\}$, describe the circuit that prepares it from
$\ket{0\cdots0}$ using only $X$ gates.
\begin{enumerate}[label=(\alph*)]
\item What is the total gate count?
\item What is the circuit \emph{depth} (can the $X$ gates all be applied
in parallel, or must they be sequential)? Contrast this with the general
amplitude state-preparation problem discussed below.
\item By the way ... how do you get the $\ket{0\cdots0}$ to begin with?
\end{enumerate}
(See the accompanying code exercise ... you'll do a simple example)


%%%%%%%%%%%%%%%%%%%%%%%%%%%%%%%%%%%%
%
%  General state preparation by rotations
%
%%%%%%%%%%%%%%%%%%%%%%%%%%%%%%%%%%%%

\item \textbf{General state preparation by rotations.}

Consider the circuit
\[ \Qcircuit @C=1em @R=.8em { \lstick{|0\rangle} & \gate{R_Y(\alpha)} & \ctrlo{1} & \ctrl{1} & \qw \\ \lstick{|0\rangle} & \qw & \gate{R_Y(\beta_0)} & \gate{R_Y(\beta_1)} & \qw } \]

where the second gate is a uniformly controlled rotation:
if the control qubit is in state $|0\rangle$, apply
$R_Y(\beta_0)$; if the control qubit is in state $|1\rangle$,
apply $R_Y(\beta_1)$.

\begin{enumerate}[label=(\alph*)]

\item Show that the final state is

\[
\begin{aligned}
|\psi\rangle
=&
\cos\frac{\alpha}{2}\cos\frac{\beta_0}{2}\,|00\rangle
+
\cos\frac{\alpha}{2}\sin\frac{\beta_0}{2}\,|01\rangle \\
&
+\sin\frac{\alpha}{2}\cos\frac{\beta_1}{2}\,|10\rangle
+
\sin\frac{\alpha}{2}\sin\frac{\beta_1}{2}\,|11\rangle .
\end{aligned}
\]

\item Consider a general real two-qubit state

\[
|\psi\rangle
=
a_{00}|00\rangle
+a_{01}|01\rangle
+a_{10}|10\rangle
+a_{11}|11\rangle ,
\]

with

\[
a_{00}^2+a_{01}^2+a_{10}^2+a_{11}^2=1.
\]

Show that it can be written in the form

\[
|\psi\rangle
=
r_0 |0\rangle |\psi_0\rangle
+
r_1 |1\rangle |\psi_1\rangle ,
\]

where $|\psi_0\rangle$ and $|\psi_1\rangle$ are normalized
one-qubit states

\[
|\psi_0\rangle
=
\frac{a_{00}|0\rangle+a_{01}|1\rangle}{r_0},
\qquad
|\psi_1\rangle
=
\frac{a_{10}|0\rangle+a_{11}|1\rangle}{r_1}.
\]

and

\[
r_0=\sqrt{a_{00}^2+a_{01}^2},
\qquad
r_1=\sqrt{a_{10}^2+a_{11}^2}.
\]

\item The first rotation $R_Y(\alpha)$ should create the state

\[
r_0|0\rangle+r_1|1\rangle .
\]

Since

\[
R_Y(\alpha)|0\rangle
=
\cos\frac{\alpha}{2}|0\rangle
+
\sin\frac{\alpha}{2}|1\rangle ,
\]

it follows that

\[
\cos\frac{\alpha}{2}=r_0,
\qquad
\sin\frac{\alpha}{2}=r_1.
\]

\begin{enumerate}[label=(\roman*)]

\item Find an expression for $\alpha$.

\item Determine formulas for $\beta_0$ and $\beta_1$.

\end{enumerate}

\item Use the above method to determine
$\alpha$, $\beta_0$, and $\beta_1$ for

\[
|\phi\rangle
=
\frac12|00\rangle
+\frac12|01\rangle
+\frac1{\sqrt8}|10\rangle
+\sqrt{\frac38}|11\rangle .
\]

\item Explain how this construction suggests a recursive procedure
for preparing an arbitrary real $n$-qubit state.

\end{enumerate}

%%%%%%%%%%%%%%%%%%%%%%%%%%%%%%%%%%%%
%
% Counting the cost of a controlled-rotation state-prep tree
%
%%%%%%%%%%%%%%%%%%%%%%%%%%%%%%%%%%%%

\item \textbf{Counting the cost of a controlled-rotation state-prep
tree.}

A uniformly controlled $R_Y$ rotation on one target qubit with
$k$ control qubits associates a rotation angle

\[
\theta_x,
\qquad
x\in\{0,1\}^k,
\]

with each computational basis state of the control register.

For example, when $k=2$, the gate is specified by four angles

\[
\theta_{00},\qquad
\theta_{01},\qquad
\theta_{10},\qquad
\theta_{11}.
\]

\begin{enumerate}[label=(\alph*)]

\item How many independent rotation angles are required for a
uniformly controlled $R_Y$ gate with $k$ control qubits?

\item Explain why a cost proportional to $2^k$ is plausible.

\textit{Hint:} The gate contains $2^k$ independent continuous
parameters.

\item Assume that a uniformly controlled $R_Y$ gate with $k$
controls can be implemented using at most

\[
c\,2^k
\]

elementary gates, where $c$ is a constant independent of $k$.

Consider the recursive state-preparation procedure from the previous
exercise: one uniformly controlled $R_Y$ layer is used to separate one
qubit from an arbitrary real $n$-qubit state, after which the procedure
recurses on the two resulting $(n-1)$-qubit branches.

Let $T(n)$ denote the number of elementary gates required.

Show that

\[
T(n)
=
2T(n-1)
+
c\,2^{n-1},
\qquad
T(1)=O(1).
\]

Clearly identify the origin of each term.

\item Divide the recurrence by $2^n$ and show that

\[
\frac{T(n)}{2^n}
=
\frac{T(n-1)}{2^{n-1}}
+
\frac{c}{2}.
\]

Hence solve the recurrence and show that

\[
T(n)=\Theta(n\,2^n).
\]

\item Compare this result with the frequently quoted upper bound
$O(4^n)$.

Why is $O(4^n)$ not a tight bound for this state-preparation procedure?
Where might a scaling proportional to $4^n$ arise naturally instead?

\end{enumerate}

%%%%%%%%%%%%%%%%%%%%%%%%%%%%%%%%%%%%
%
%  Complex amplitudes
%
%%%%%%%%%%%%%%%%%%%%%%%%%%%%%%%%%%%%

\item \textbf{Complex amplitudes.}
Repeat Exercise 2, but now allow \emph{complex} amplitudes $a_i \in
\mathbb C$ (so each ``peeling off a qubit'' step needs both an $R_Y$
magnitude rotation \emph{and} an $R_Z$ phase rotation per branch). Does
the asymptotic growth rate (as a function of $n$) change from the
real-amplitude case, or only the constant factor?

%%%%%%%%%%%%%%%%%%%%%%%%%%%%%%%%%%%%
%
%  QRAM and the ``log-depth mirage''
%
%%%%%%%%%%%%%%%%%%%%%%%%%%%%%%%%%%%%

\item \textbf{QRAM and the ``log-depth mirage''.}
A common claim is: ``QRAM lets you load $N$ data points in $\mathcal
O(\log N)$ time.'' The bucket-brigade QRAM architecture routes an
address query $\ket i$ down a binary tree of $\mathcal O(N)$ routing
qubits to reach memory cell $i$ in circuit \emph{depth} $\mathcal
O(\log N)$.
\begin{enumerate}[label=(\alph*)]
\item Suppose the $N$ values already sit in the QRAM's memory cells as
fixed physical parameters (like a lookup table burned into hardware).
Explain why the query
$\ket i\ket 0 \mapsto \ket i \ket{x_i}$
can indeed be performed in $\mathcal O(\log N)$ depth.
\item Now suppose instead that you want to encode $N$ \emph{arbitrary}
real amplitudes $a_0,\dots,a_{N-1}$ (freshly computed classically, not
sitting in a pre-built table) into the state $\sum_i a_i \ket i$.
Explain why this "loading" step still requires performing $N$
\emph{genuinely different} single-qubit rotations $R_Y(\theta_i)$ --
one per memory cell -- and therefore costs $\Omega(N)$ gates overall,
regardless of the $\mathcal O(\log N)$ query \emph{depth} of the
underlying QRAM routing structure.
\item Reconcile (a) and (b): under what circumstances does QRAM give a
genuine complexity advantage, and when is the "$\mathcal O(\log N)$"
figure misleading about the \emph{total} cost of getting your data into
the machine?
\end{enumerate}

%%%%%%%%%%%%%%%%%%%%%%%%%%%%%%%%%%%%
%
%  Structured states are cheap
%
%%%%%%%%%%%%%%%%%%%%%%%%%%%%%%%%%%%%

\item \textbf{Structured states are cheap.}
Give an example of an $n$-qubit state with (generically) all $2^n$
amplitudes nonzero, that can nonetheless be prepared in $\mathrm{poly}(n)$
gates. (Hint: consider the state produced by applying the Quantum
Fourier Transform to a computational basis state, or a simple product
state written in a different basis, or the GHZ or $W$ states.) Contrast
this with a generic (Haar-random) state, which Exercises 2--3 show needs
exponentially many gates. What structural property distinguishes the
cheap states from the generic ones?

\end{enumerate}

\end{document}
